% Generated draft: requires mathematical and source review.
\documentclass[11pt]{article}
\usepackage[margin=25mm]{geometry}
\usepackage{fontspec}
\setmainfont{Latin Modern Roman}
\usepackage{amsmath,amssymb,mathtools,unicode-math}
\setmathfont{Latin Modern Math}
\usepackage{enumitem,needspace,xurl,hyperref}
\setlength{\parindent}{0pt}
\setlength{\parskip}{5pt}
\setlength{\emergencystretch}{3em}
\newcommand{\N}{\mathbb{N}}
\newcommand{\Z}{\mathbb{Z}}
\newcommand{\Q}{\mathbb{Q}}
\newcommand{\R}{\mathbb{R}}
\newcommand{\C}{\mathbb{C}}
\newcommand{\F}{\mathbb{F}}
\newcommand{\A}{\mathbb{A}}
\newcommand{\PP}{\mathbb{P}}
\newcommand{\E}{\mathbb{E}}
\newcommand{\eps}{\varepsilon}
\newcommand{\dd}{\,\mathrm{d}}
\begin{document}
\section*{UnsolvedMath ID: 30004612}
\textbf{Source problem code:} catalogue entry
\subsection{Definitions and mathematical statement}
An abelian surface over \(\Q\) is an abelian variety of dimension \(2\) defined over \(\Q\). Two such surfaces are \(\Q\)-isogenous if there is a surjective homomorphism between them, defined over \(\Q\), with finite kernel. For a positive integer \(N\), let
\[
\mathcal A_N=\left\{[A]_{\Q\text{-isog}}: A/\Q\text{ is an abelian surface of conductor }N\right\},
\]
where the conductor is the usual arithmetic conductor of the abelian variety.

The supplied question asks whether the Brumer--Kramer paramodular correspondence gives a bijection between \(\mathcal A_N\) and the appropriate paramodular forms of level \(N\). If \(\mathcal P_N\) denotes the intended collection of forms, the requested assertion has the schematic form
\[
\Phi_N:\mathcal A_N\xrightarrow{\;\sim\;}\mathcal P_N,
\]
where \(\Phi_N\) must be the intended Brumer--Kramer correspondence.

\textbf{Unresolved formulation.} The source does not define \(\mathcal P_N\), specify an equivalence relation on its forms, or characterize \(\Phi_N\). It supplies no weight, eigenform or newform conditions, coefficient-field conditions, normalization, or treatment of lifts. It also gives no restrictions on the abelian surfaces beyond their dimension, field of definition, and conductor. Thus \(\mathcal A_N\) above records the literal source wording; it is not a verified identification of the domain of the precise Brumer--Kramer conjecture. No restriction on \(N\) is stated, although one supplied reference concerns odd conductor. The displayed bijection records the intended mathematical target, but the missing specifications prevent a complete, unambiguous conjectural statement from being recovered from this record alone.
\subsection{Short English statement}
The question asks whether the Brumer--Kramer correspondence pairs the relevant \(\Q\)-isogeny classes of abelian surfaces of conductor \(N\) bijectively with the relevant paramodular forms of level \(N\). The supplied record leaves the required restrictions on surfaces, forms, and levels unspecified, so this is a qualified statement-only draft requiring mathematical review.
\subsection{Sources}
\begin{itemize}
\item Local UnsolvedMath registry, permanent ID \(30004612\), title \emph{Paramodular Forms and Isogeny Classes of Abelian Surfaces}, source problem code \texttt{catalogue entry}, in the Oberwolfach Reports Open Problems collection. The record attributes the conjecture to Brumer--Kramer, lists proposed year \(2021\), and names no individual proposer.
\item Supplied report citation: \emph{Moduli spaces and Modular forms (2022)}, Oberwolfach Reports, pp.~301--359, DOI \texttt{10.4171/owr/2021/6}. \url{https://doi.org/10.4171/owr/2021/6}. The background describes these reports as records of mathematical workshops and their open-problem discussions. The supplied literature list labels this same reference \emph{Oberwolfach Report 6/2021}.
\item Brumer--Kramer, \emph{Paramodular Abelian Varieties of Odd Conductor}. \url{https://arxiv.org/abs/1004.4699}.
\item MIT lecture notes, supplied label: \emph{2024 status summary}. \url{https://math.mit.edu/classes/18.786/2024/LectureNotes23.pdf}.
\item Supplied literature assessment, checked \texttt{2026-08-22}: the correspondence has many verified examples and computational evidence, but both directions of the general conjecture remain open. This is supplied dated curation derived from the OWR record, not an independently verified literature assessment. The supplied status is \texttt{open}; the assessment and cited claims require independent verification.
\end{itemize}
\end{document}
